\section{Introduction}
\label{sec:introduction}

\IEEEPARstart{A}{utonomous} systems increasingly rely on runtime layers that
monitor safety, performance, or mission properties and intervene during
execution.  Examples include modular runtime assurance for robotics
\cite{desai2019soter}, mission-level runtime assurance for autonomous driving
\cite{tsai2026missionrta}, compositional runtime verification for robot swarms
\cite{kekatos2026swarmrta}, and predictive runtime verification, enforcement,
and monitoring
\cite{pinisetty2017prv,pinisetty2017pre,yoon2019prm,wheatman2021radics}.
Prediction is attractive because the extra lead it provides may allow a recovery
that preserves the mission instead of an emergency stop.

A predictive monitor is, however, an inference component inside the robot's
decision cycle.  It must estimate how fast a margin is degrading from noisy and
possibly delayed observations, run within the cycle's latency budget, and be
tuned and certified like any other safety function.  A reactive rule can also
intervene early: it only needs a larger fixed margin.  Comparing a predictor
with an unpadded reactive trigger therefore does not show that prediction is
necessary.  The fair comparison gives the fixed margin a comparable tuning
budget.  We ask: \emph{when does predicting the loss of feasibility outperform
reacting to it?}

\Cref{sec:analysis} answers this question within an idealized timing model.  A
trigger that fires when the projected time to margin expiry falls below a
horizon is a reactive trigger with a rate-scaled margin.  An earlier trigger
changes a recovery outcome only within a band set by recovery duration and
deadline slack, and observation age shrinks that band exactly as lost lead does.
The minimal predictive horizon wastes less lead than the minimal safe fixed pad
when recovery requires a similar lead in every case but, because degradation
rates differ, a different margin.  \Cref{sec:method,sec:results} examine these
conditions in controlled simulations.  Exploratory sweeps locate the boundary,
and two further evaluations separate tuning from testing.  In the held-out
single-margin comparison, the tuned fixed rule records $0/6750$ conflict
collisions and the tuned predictive rule records $2/6750$, so the outcomes do not
support an efficiency comparison.  In a three-robot study with genuinely
conflicting geometry and an explicit serialization state machine, both tuned
rules succeed in every held-out conflict episode, and the predictive rule has a
shorter mean makespan on the fixed scenario grid.

The paper makes four contributions:
\begin{enumerate}
  \item a reduction showing that the studied horizon-based predictive trigger is
    a rate-adaptive current-margin threshold (\cref{prop:reduction});
  \item a boundary condition, in terms of recovery duration, deadline slack, and
    observation age, for when added lead changes the recovery outcome
    (\cref{prop:boundary,cor:age}), and a characterization of when the minimal
    predictive horizon wastes less lead than the minimal safe fixed pad
    (\cref{prop:homogeneity});
  \item a reproducible evaluation protocol that separates exploratory in-sample
    sweeps from seed-disjoint tuning and testing, gives both rules equal tuning
    budgets, reports failures before timing, and withholds timing inference when
    failure outcomes differ; and
  \item a genuine-conflict three-robot evaluation with explicit approaching,
    waiting, crossing, and cleared states, which identifies a one-shot
    serialization regime in which the rate-adaptive trigger shortens completion
    time relative to the tuned fixed pad at matched observed outcomes.
\end{enumerate}

The paper characterizes a trigger mechanism rather than proposing a controller.
The experiments are stylized two-dimensional kinematic simulations chosen to
isolate that mechanism, and the results apply only to the tested models.
\Cref{sec:limitations} details these limits.
